\documentclass[a4paper]{article}
%\usepackage{simplemargins}

%\usepackage[square]{natbib}
\usepackage{amsmath}
\usepackage{amsfonts}
\usepackage{amssymb}
\usepackage{graphicx}
\usepackage{booktabs}
\usepackage{adjustbox}
\usepackage{xcolor}
\usepackage{placeins}

\usepackage[
  top=1in,
  bottom=1in,
  left=1.25in,
  right=1.25in
]{geometry}

\begin{document}
\pagenumbering{gobble}

\Large
 \begin{center}
A multimodal spatial-attention multiple-instance learning method for cervical cancer recurrence prediction on pretreatment planning CT and PET imaging data\\ 

\hspace{10pt}

% Author names and affiliations
\large
Jialin Yue$^{a}$, Zong Fan$^{b}$, Hua Li$^{a,b}$*\\
\normalsize
$^{a}$ Department of Radiation Oncology, Washington University in St. Louis, St. Louis, MO, United States\\
$^{b}$ Department of Bioengineering, University of Illinois Urbana-Champaign, Urbana, IL, United States\\
\footnotesize *Corresponding author: hua.li@wustl.edu

\hspace{10pt}



\end{center}

\hspace{10pt}

\textbf{Introduction:} Accurate prediction of cervical cancer recurrence using pre-radiotherapy imaging is critical for optimizing treatment planning and risk stratification to facilitate personalized therapeutic strategies. 
While integrating multimodal data, such as planning CT and PET, significantly enhances predictive power over single-modality approaches, it introduces substantial technical challenges. 
Specifically, multimodal datasets often exhibit protocol variability across clinical sites, subtle prognostic features that are easily overlooked, and missing modalities where patients may lack either the PET or CT component.
To address these limitations, multimodal AI-driven approaches, particularly Multiple-Instance Learning (MIL), offer a promising methodology for robust local information aggregation. 
In this work, we propose a tumor-neighborhood spatial-attention MIL approach designed to improve outcome prediction by processing multiple overlapping subvolumes from the tumor and its neighboring area. 
This methodology reduces sensitivity to local nuisance variations and modality incompleteness while preserving the essential information provided by PET/CT. 
Furthermore, our model offers inherent interpretability through learned attention weights, allowing for the evidence-based identification of specific tumor-neighborhood areas that most significantly contribute to the patient-level recurrence prediction.

%local infomation/knowledge
\textbf{Method:}
We propose a tumor-neighborhood spatial-attention MIL model that characterizes the tumor and its surrounding environment through the extraction of multiple 3D subvolumes from planning CT and PET. 
The model architecture consists of: (1) dual modality-specific DLA3D (Deep Layer Aggregation 3D) encoders optimized to learn high-dimensional local feature representations from CT and PET subvolumes, 
and (2) a gated-attention MIL pooling module that dynamically aggregates information across multiple tumor-neighborhood regions to generate a global feature for patient-level prediction, thus improving overall prediction accuracy. 
Additionally, this modular feature extraction and flexible fusion architecture enable the model to process available input sources effectively while maintaining robustness in the presence of missing modalities.

To enable model interpretability for transparency in medical AI, the gated-attention mechanism assigns localized importance scores to each tumor-neighborhood subvolume, 
representing its relative contribution to the final recurrence prediction. 
This way, these weights are utilized to identify highly informative regions within the multimodal space. 
We validate the clinical relevance of these attention-derived importance through an ablation study, where the top-ranked subvolumes (Top-3) are selectively removed at inference. 
By re-evaluating model performance post-removal, we quantify the performance drop in predictive probability (Fig.~\ref{fig:top3_removal}); 
a substantial decrease in confidence indicates that the attention mechanism has successfully localized the most significant prognostic cues within the tumor-neighborhood.


\textbf{Results:}
 patient dataset. In our work, we frame recurrence outcome prediction as a binary classification task (recurrent vs non-recurrent). Using five-fold cross-validation, mean AUC improved from 0.672 to 0.713 (+0.041), and mean accuracy improved from 0.572 to 0.666 (+0.094). Specificity increased from 0.485 to 0.696 (+0.211), while sensitivity decreased from 0.715 to 0.627 (-0.088). Mean FP decreased from 15.2 to 9.2 (-6.0), and mean FN increased from 4.8 to 6.2 (+1.4) (Table~\ref{tab:mil_cv_results}). For (what?)interpretability, we conducted a Top-3 patch removal test on two folds: fold1 with the highest AUC showed marginal significance(what for) ($p=0.051$), and fold4 with the largest overall metric gain showed statistical significance ($p<0.05$) (Fig.~\ref{fig:top3_removal}).

\textbf{Conclusions:}
Our approach incorporates a registered planning CT + PET baseline with a lightweight MIL attention pooling module to reduce reliance on a single tumor-centered view and improve robustness under lesion heterogeneity and structured nuisance signals. By aggregating evidence from multiple tumor-neighborhood subvolumes, the model achieves improved outcome prediction performance. In addition, the attention pooling weights provide instance-level interpretability by ranking informative patches, and this evidence ranking is supported by our top-ranked patch removal analysis.

\textbf{Potential impacts:}
Our method provides a lightweight path toward more robust pelvic PET/CT outcome prediction, with potential extensions to cross-site generalization, richer region design and weakly supervised ROI learning, and controlled fusion of additional modalities. It also offers a new interpretability angle for PET/CT through instance-level evidence ranking.

% in preamble:
% \usepackage{booktabs}
% \usepackage{adjustbox}

\FloatBarrier
\newpage
\thispagestyle{plain}

\begin{center}
{\large\bfseries Supplementary Page}
\end{center}
\vspace{-0.2em}

\FloatBarrier  % 防止下面的 floats 跑到标题上面


\begin{table}[!h]
\centering
\footnotesize
\setlength{\tabcolsep}{2.8pt}
\renewcommand{\arraystretch}{1.12}

\begin{adjustbox}{max width=\linewidth}
\begin{tabular}{l *{16}{c}}
\toprule
 & \multicolumn{2}{c}{AUC}
 & \multicolumn{2}{c}{TP}
 & \multicolumn{2}{c}{TN}
 & \multicolumn{2}{c}{FP}
 & \multicolumn{2}{c}{FN}
 & \multicolumn{2}{c}{Sensitivity}
 & \multicolumn{2}{c}{Specificity}
 & \multicolumn{2}{c}{Accuracy} \\
\cmidrule(lr){2-3}\cmidrule(lr){4-5}\cmidrule(lr){6-7}\cmidrule(lr){8-9}
\cmidrule(lr){10-11}\cmidrule(lr){12-13}\cmidrule(lr){14-15}\cmidrule(lr){16-17}
Fold & w/o & w/ & w/o & w/ & w/o & w/ & w/o & w/ & w/o & w/ & w/o & w/ & w/o & w/ & w/o & w/ \\
\midrule
1 & 0.739 & \textbf{0.772}
  & 10 & 9
  & 20 & 19
  & 11 & 12
  & 2 & 3
  & \textbf{0.833} & 0.750
  & \textbf{0.645} & 0.613
  & \textbf{0.698} & 0.651 \\
2 & 0.666 & \textbf{0.763}
  & 17 & 13
  & 8 & 20
  & 18 & 6
  & 2 & 6
  & \textbf{0.895} & 0.684
  & 0.308 & \textbf{0.769}
  & 0.556 & \textbf{0.733} \\
3 & \textbf{0.734} & 0.704
  & 10 & 5
  & 17 & 22
  & 10 & 5
  & 6 & 11
  & \textbf{0.625} & 0.313
  & 0.630 & \textbf{0.815}
  & \textbf{0.628} & 0.628 \\
4 & 0.595 & \textbf{0.658}
  & 10 & 12
  & 13 & 22
  & 18 & 9
  & 8 & 6
  & 0.556 & \textbf{0.667}
  & 0.419 & \textbf{0.710}
  & 0.469 & \textbf{0.694} \\
5 & 0.626 & \textbf{0.670}
  & 12 & 13
  & 14 & 19
  & 19 & 14
  & 6 & 5
  & 0.667 & \textbf{0.722}
  & 0.424 & \textbf{0.576}
  & 0.510 & \textbf{0.627} \\
\midrule
Mean & 0.672 & \textcolor{blue}{\textbf{0.713}}
     & \multicolumn{2}{c}{--}
     & \multicolumn{2}{c}{--}
     & \multicolumn{2}{c}{--}
     & \multicolumn{2}{c}{--}
     & \textcolor{blue}{\textbf{0.715}} & 0.627
     & 0.485 & \textcolor{blue}{\textbf{0.696}}
     & 0.572 & \textcolor{blue}{\textbf{0.666}} \\
\bottomrule
\end{tabular}
\end{adjustbox}
\caption{Five-fold performance comparison between models trained without and with MIL. For AUC, sensitivity, specificity, and accuracy, the better value within each fold is highlighted in bold; the better mean value is highlighted in blue. Confusion-matrix counts are omitted for the mean row.}
\label{tab:mil_cv_results}
\end{table}
\vspace{-1.0em}


\begin{figure}[!h]
  \centering
  \includegraphics[width=\linewidth]{fold1_fold4_top3remove.pdf}
  \caption{Ablation-based interpretability via Top-3 patch removal. Left: the fold achieving the highest AUC (Mann--Whitney $p=0.051$). Right: the fold with the largest overall metric gain (Mann--Whitney $p<0.001$).
For each patient, we rank tumor-neighborhood subvolumes by MIL attention weights and remove the Top-3 ranked patches at inference.
The y-axis shows the log-odds drop, $\Delta=\mathrm{logit}(p_{\text{base}})-\mathrm{logit}(p_{\text{removed}})$, measuring the decrease in predicted probability after removal.
The x-axis compares cases whose removed Top-3 patches have high versus low within-patch intensity standard deviation (STD), defined by tertiles across patients (top vs bottom tertile).}
 
  \label{fig:top3_removal}
\end{figure}



\end{document}